\chapter*{\Large \center Abstract}

Prompt-based Large Language Models (LLMs), such as GPT-style systems, Gemini, and LLaMA, demonstrate strong zero-shot performance in Grammatical Error Correction (GEC). However, their generative fluency can lead to over-correction, whereby models introduce unnecessary lexical or structural changes beyond minimal grammatical correction. For second-language (L2) learners, such rewriting may obscure the source of errors and reduce the pedagogical value of feedback\cite{loem2023exploringeffectivenessgpt3grammatical}.

This dissertation investigates how prompt engineering influences over-correction in LLM-based GEC. It compares multiple prompting strategies and evaluates model outputs using ERRANT-based accuracy metrics alongside a Composite Over-Correction Index (OCI), which combines source-output divergence with signals of fluency improvement and meaning preservation to provide an indicator of over-correction risk.

The findings suggest that prompt design plays an important role in how LLMs correct learner writing. More constrained prompts can reduce unnecessary changes while still maintaining a similar level of grammatical accuracy. Overall, this study provides a framework for analysing over-correction and offers guidance for designing LLM-based writing support tools that make minimal but useful corrections. The aim is to support L2 learners’ writing development without replacing their own voice\cite{park-etal-2025-leveraging}.
